% ROZDZIAŁ 18 — Radon-Nikodym, pełny dowód przez przestrzeń Hilberta
\chapter[Twierdzenie Radona--Nikodyma]{Twierdzenie Radona--Nikodyma i pochodne miar}
\label{ch:radon-nikodym-18}

\providecommand{\atlasref}[2]{#2}


\section{Problem reprezentacji miary przez funkcję}
Niech $(X,\mathcal F,\mu)$ będzie przestrzenią z miarą,
a $\nu$ drugą miarą nieujemną na tej samej $\sigma$-algebrze.
Chcielibyśmy znaleźć mierzalną funkcję $f\ge0$, dla której
\begin{equation}\label{eq:RN-question-18}
\nu(E)=\int_E f(x)\,\mathrm d\mu(x)
\qquad(E\in\mathcal F).
\end{equation}
Wynikałoby stąd natychmiast $\nu\ll\mu$.
Czy warunek ten jest wystarczający? Odpowiedź jest twierdząca
przy odpowiednich założeniach o skończoności miar.

\begin{definicja}[Pochodna Radona--Nikodyma]
\label{def:RN-derivative-18}
Jeśli $\nu$ i $\mu$ są miarami nieujemnymi na $(X,\mathcal F)$
i istnieje mierzalna funkcja $f:X\to[0,\infty]$
spełniająca \eqref{eq:RN-question-18}, nazywamy ją
\emph{gęstością} lub \emph{pochodną Radona--Nikodyma}
miary $\nu$ względem $\mu$ i piszemy
$f=\mathrm d\nu/\mathrm d\mu$. Jeżeli taka pochodna
jest jednoznaczna z dokładnością do równości $\mu$-prawie
wszędzie, symbol ten oznacza jej klasę równoważności.
W tej definicji symbol $\mathrm d\nu/\mathrm d\mu$
nie jest zwykłym ilorazem dwóch funkcji.
\end{definicja}

\begin{lemat}[Całkowanie względem miary o gęstości]
\label{lem:integration-density-18}
Niech $\lambda$ będzie miarą nieujemną na $(X,\mathcal F)$,
$w:X\to[0,\infty]$ funkcją mierzalną. Wówczas wzór
$\mu(E)=\int_Ew\,\mathrm d\lambda$ dla $E\in\mathcal F$
określa miarę nieujemną na $(X,\mathcal F)$.
Ponadto dla każdej mierzalnej $\varphi:X\to[0,\infty]$
zachodzi
\begin{equation}\label{eq:integration-density-18}
\int_X\varphi(x)\,\mathrm d\mu(x)
=\int_X\varphi(x)w(x)\,\mathrm d\lambda(x),
\end{equation}
przy konwencji $0\cdot\infty=0$ dla iloczynu
nieujemnych funkcji mierzalnych. Jeśli $\varphi:X\to\mathbb C$ jest mierzalna, to
$\varphi\inL^1(\mu)$ wtedy i tylko wtedy, gdy
\[
 \int_X|\varphi|w\,\mathrm d\lambda<\infty.
\]
W tym przypadku analogiczna równość zachodzi dla całek
zespolonych. W iloczynie $\varphi w$ przyjmujemy wartość
zero tam, gdzie $\varphi=0$, także gdy $w=+\infty$;
na ewentualnym $\lambda$-zerowym zbiorze
$\{\varphi\ne0,\ w=+\infty\}$ wybieramy wartość zero
jako skończonego reprezentanta iloczynu.
\end{lemat}
\begin{proof}
Najpierw sprawdzamy, że $\mu$ jest miarą. Całka
po zbiorze pustym wynosi zero. Jeżeli $(A_j)_{j\ge1}$
jest rodziną parami rozłącznych zbiorów mierzalnych, to
przy konwencji $0\cdot\infty=0$ mamy
\[
 \sum_{j=1}^{N}w\chi_{A_j}
 \uparrow w\chi_{\bigcup_{j\ge1}A_j}.
\]
Z addytywności całki nieujemnej dla skończonych sum
oraz z twierdzenia \atlasref{ME-08::thm:mct}{ME-08: twierdzenie Beppo Leviego} o zbieżności
monotonicznej otrzymujemy
\[
 \mu\Bigl(\bigcup_{j\ge1}A_j\Bigr)
 =\lim_{N\to\infty}\sum_{j=1}^{N}
          \int_X w\chi_{A_j}\,\mathrm d\lambda
 =\sum_{j\ge1}\mu(A_j).
\]
Dowód nie wymaga $\sigma$-skończoności $\lambda$
ani skończoności $w$. Dla $\varphi=\chi_A$
tożsamość całkowa jest definicją $\mu$.
Z liniowości dla funkcji prostych nieujemnych otrzymujemy
równość dla wszystkich takich funkcji. Każdą funkcję
$\varphi\ge0$ aproksymujemy z dołu ciągiem mierzalnych
funkcji prostych $\varphi_n\uparrow\varphi$.
Dwukrotne zastosowanie twierdzenia \atlasref{ME-08::thm:mct}{ME-08: twierdzenie Beppo Leviego} o zbieżności monotonicznej
daje \eqref{eq:integration-density-18}. W przypadku
funkcji zespolonej stosujemy tę tożsamość najpierw do $|\varphi|$.
Otrzymujemy równość całek modułu, która daje obie
implikacje kryterium całkowalności. Jeżeli są skończone,
to mierzalny zbiór $\{\varphi\ne0,\ w=+\infty\}$ jest
$\lambda$-zerowy, ponieważ $|\varphi|w=+\infty$ na nim.
Zatem opisany w tezie skończony reprezentant iloczynu
jest mierzalny i całkowalny. Następnie stosujemy równość do części
dodatnich i ujemnych części rzeczywistej i urojonej $\varphi$,
które są całkowalne względem $\mu$ dzięki założeniu, oraz łączymy cztery równości.
\end{proof}

\section{Krok pierwszy: miary skończone}
\begin{twierdzenie}[Radon--Nikodym dla miar skończonych]
\label{thm:RN-finite-18}
Niech $\mu,\nu$ będą skończonymi miarami nieujemnymi
na $(X,\mathcal F)$ i niech $\nu\ll\mu$.
Istnieje mierzalna funkcja $f:X\to[0,\infty)$ taka, że
\begin{equation}\label{eq:RN-finite-18}
\nu(E)=\int_E f(x)\,\mathrm d\mu(x)
\quad\text{dla każdego }E\in\mathcal F.
\end{equation}
Funkcja $f$ należy do $L^1(\mu)$, jest jednoznaczna
$\mu$-prawie wszędzie i
$\int_Xf\,\mathrm d\mu(x)=\nu(X)$.
\end{twierdzenie}
\begin{proof}
\emph{Etap 1: wspólna miara dominująca.}
Definiujemy $\lambda=\mu+\nu$. Jest to miara skończona
na $\mathcal F$. Dla $u\in L^2(\lambda;\mathbb C)$
kładziemy
\[
T(u)=\int_X u(x)\,\mathrm d\nu(x).
\]
Funkcjonał ten jest dobrze określony: ponieważ
$\nu\le\lambda$ w sensie wartości na zbiorach,
$\int_X|u|^2\,\mathrm d\nu(x)\le
\int_X|u|^2\,\mathrm d\lambda$.
Nierówność Cauchy'ego--Schwarza daje
\[
|T(u)|\le\nu(X)^{1/2}
\left(\int_X|u|^2\,\mathrm d\nu(x)\right)^{1/2}
\le\nu(X)^{1/2}\|u\|_{L^2(\lambda)}.
\]
Jeśli dwa przedstawiciele $u$ są równe $\lambda$-prawie
wszędzie, to są równe $\nu$-prawie wszędzie, więc
$T$ jest funkcjonałem na klasach $L^2(\lambda)$.

\emph{Etap 2: reprezentacja Hilbertowska.}
Wykorzystujemy twierdzenie Riesza o reprezentacji
funkcjonałów na przestrzeni Hilberta $L^2(\lambda)$:
istnieje $h\in L^2(\lambda)$ takie, że
\begin{equation}\label{eq:RN-lambda-h-18}
T(u)=\int_Xu(x)h(x)\,\mathrm d\lambda(x)
\qquad(u\in L^2(\lambda)).
\end{equation}
Dla kompletności przypomnijmy uzasadnienie użytej
reprezentacji. Jeśli $T=0$, bierzemy $h=0$.
W przeciwnym razie $K=\ker T$ jest domkniętą
podprzestrzenią Hilberta; każdy $v$ ma jednoznaczny
rzut ortogonalny $v_K$ na $K$, którego istnienie wynika
z zupełności i tożsamości równoległoboku: ciąg
minimalizujący odległość od $K$ jest ciągiem Cauchy'ego.
Wybieramy $v\in L^2(\lambda)$ z $T(v)=1$ i rzutujemy
$v$ na $K^\perp$. Otrzymany wektor $u$ spełnia
$u\in K^\perp$ i $T(u)=1$, ponieważ $T$ zeruje się na $K$.
Dla dowolnego $v$ różnica $v-T(v)u$ należy do $K$,
więc przy iloczynie skalarnym liniowym w pierwszym
argumencie mamy
$\langle v,u\rangle=T(v)\|u\|_2^2$.
Zatem $T(v)=\langle v,u\rangle/\|u\|_2^2$,
a reprezentantem w zapisie \eqref{eq:RN-lambda-h-18}
jest $h=\overline u/\|u\|_2^2$.
Nie użyliśmy w tym miejscu twierdzenia o gęstościach miar.

\emph{Etap 3: oszacowanie $0\le h\le1$.}
Ponieważ $\lambda(X)<\infty$, funkcje charakterystyczne $\chi_E$
należą do $L^2(\lambda)$. Z \eqref{eq:RN-lambda-h-18}
otrzymujemy
\begin{equation}\label{eq:RN-mu-nu-h-18}
\nu(E)=\int_Eh\,\mathrm d\lambda,\qquad
\mu(E)=\int_E(1-h)\,\mathrm d\lambda.
\end{equation}
Pierwsza całka jest rzeczywista i nieujemna dla
każdego $E$. Wybierając kolejno zbiory, na których
$\operatorname{Im}h$ jest dodatnia lub ujemna,
dostajemy $\operatorname{Im}h=0$ $\lambda$-prawie wszędzie.
Podobnie z nieujemności obu całek w
\eqref{eq:RN-mu-nu-h-18} otrzymujemy
$h\ge0$ i $1-h\ge0$ prawie wszędzie.
Po modyfikacji na zbiorze $\lambda$-zerowym możemy
więc założyć, że $0\le h\le1$ wszędzie.

\emph{Etap 4: wykluczenie obszaru $h=1$.}
Niech $S=\{x:h(x)=1\}$. Druga równość
w \eqref{eq:RN-mu-nu-h-18} daje $\mu(S)=0$.
Z założenia $\nu\ll\mu$ wynika $\nu(S)=0$.
Zatem $\lambda(S)=\mu(S)+\nu(S)=0$.
Poza $S$ mianownik $1-h$ jest dodatni.
Definiujemy
\[
f(x)=\begin{cases}
\dfrac{h(x)}{1-h(x)},&h(x)<1,\\[4pt]
0,&h(x)=1.
\end{cases}
\]
Funkcja $f$ jest mierzalna i nieujemna. Na mocy
lematu~\ref{lem:integration-density-18}, zastosowanego
do $\mu(E)=\int_E(1-h)\,\mathrm d\lambda$,
\[
\int_E f\,\mathrm d\mu(x)
 =\int_E f(1-h)\,\mathrm d\lambda
 =\int_Eh\,\mathrm d\lambda
 =\nu(E).
\]
W środkowej równości ewentualne punkty $S$ nie
mają znaczenia, gdyż $\lambda(S)=0$.
Dla $E=X$ mamy $\int_Xf\,\mathrm d\mu(x)=\nu(X)<\infty$,
więc $f$ jest skończona $\mu$-prawie wszędzie
i należy do $L^1(\mu)$.

\emph{Etap 5: jednoznaczność.}
Jeżeli $f,g\in L^1(\mu)$, $f,g\ge0$, spełniają
\eqref{eq:RN-finite-18}, to dla każdego $E\in\mathcal F$
$\int_E(f-g)\,\mathrm d\mu(x)=0$.
Biorąc $E=\{f>g\}$, otrzymujemy
$\int_{\{f>g\}}(f-g)\,\mathrm d\mu(x)=0$.
Całkowana funkcja jest nieujemna, więc
$\mu(\{f>g\})=0$. Zamiana $f$ z $g$ daje
$\mu(\{g>f\})=0$. Stąd $f=g$ prawie wszędzie.
\end{proof}

\begin{uwaga}[Rola pomocniczej miary $\lambda=\mu+\nu$]
Dowód nie zakłada z góry, że $\nu$ ma gęstość względem
$\mu$. Najpierw reprezentujemy funkcjonał całkowy w
$L^2(\mu+\nu)$, a następnie uzyskujemy dwie gęstości
$h$ i $1-h$ względem tej samej miary $\lambda$.
Dopiero warunek $\nu\ll\mu$ pozwala wykluczyć
zbiór $\{h=1\}$ i podzielić te gęstości.
W ten sposób znika potencjalne błędne koło.
\end{uwaga}

\section{Przypadek sigma-skończony: twierdzenie główne}
\begin{twierdzenie}[Twierdzenie Radona--Nikodyma]
\label{thm:RN-sigfin-18}
Niech $\mu$ i $\nu$ będą miarami nieujemnymi,
$\sigma$-skończonymi na przestrzeni mierzalnej $(X,\mathcal F)$.
Jeśli $\nu\ll\mu$, to istnieje mierzalna funkcja
$f:X\to[0,\infty)$, skończona $\mu$-prawie wszędzie,
taka, że
\begin{equation}\label{eq:RN-general-18}
\nu(E)=\int_E f(x)\,\mathrm d\mu(x)
\qquad(E\in\mathcal F).
\end{equation}
Funkcja $f$ jest jednoznaczna $\mu$-prawie wszędzie.
Dla każdej mierzalnej funkcji $\varphi\ge0$ zachodzi
\begin{equation}\label{eq:RN-integral-change-18}
\int_X\varphi(x)\,\mathrm d\nu(x)
=\int_X\varphi(x)f(x)\,\mathrm d\mu(x).
\end{equation}
W szczególności, jeśli $\nu(X)<\infty$, to
$f\in L^1(\mu)$ i $\int_Xf\,\mathrm d\mu(x)=\nu(X)$.
\end{twierdzenie}
\begin{proof}
Ponieważ obie miary są $\sigma$-skończone,
można wybrać mierzalny, parami rozłączny podział
$X=\bigsqcup_{n\ge1}C_n$ taki, że
$\mu(C_n)<\infty$ i $\nu(C_n)<\infty$.
Jego konstrukcję podaliśmy w dowodzie twierdzenia
\ref{thm:lebesgue-decomp-sigfin-17}.
Dla ograniczeń obu miar do $C_n$ zachodzi
$\nu|_{C_n}\ll\mu|_{C_n}$. Z twierdzenia
\ref{thm:RN-finite-18} otrzymujemy mierzalną
$\mathcal F|_{C_n}$-funkcję $f_n\ge0$ taką, że
$\nu(E\cap C_n)=\int_{E\cap C_n}f_n\,\mathrm d\mu(x)$
dla wszystkich $E\in\mathcal F$.
Funkcja $f$, określona przez $f|_{C_n}=f_n$,
jest mierzalna, gdyż $C_n$ tworzą przeliczalny
podział mierzalny. Dla każdego $E$ twierdzenie
Beppo Leviego i przeliczalna addytywność $\nu$ dają
\[
\int_E f\,\mathrm d\mu(x)
=\sum_{n\ge1}\int_{E\cap C_n}f_n\,\mathrm d\mu(x)
=\sum_{n\ge1}\nu(E\cap C_n)=\nu(E).
\]
Z sumowalności $\int_{C_n}f_n\,\mathrm d\mu(x)=\nu(C_n)<\infty$
wynika, że $f$ jest skończona prawie wszędzie na każdym
$C_n$, a więc także na $X$.

Jeśli $f,g$ reprezentują $\nu$, to na każdym $C_n$
obie funkcje są całkowalne i z jednoznaczności w
przypadku skończonym są równe $\mu$-prawie wszędzie.
Po zsumowaniu przeliczalnej rodziny zbiorów wyjątkowych
otrzymujemy jednoznaczność na $X$.
Wzór \eqref{eq:RN-integral-change-18} jest zastosowaniem
lematu~\ref{lem:integration-density-18}.
\end{proof}

\begin{uwaga}[Co dokładnie zapewnia $\sigma$-skończoność?]
Założenie dotyczy \emph{obu} miar i wystarcza, aby rozbić
$X$ na przeliczalnie wiele części, na których
$\mu$ i $\nu$ są skończone. Nie jest ono warunkiem
koniecznym w każdej konkretnej przestrzeni;
istnieją szersze klasy miar, dla których reprezentacja
Radona--Nikodyma również działa. Bez dalszych warunków
nie można go jednak po prostu wykreślić.
\end{uwaga}

\begin{przyklad}[Brak $\sigma$-skończoności może zniszczyć wniosek]
\label{ex:RN-failure-nonsigma-18}
Niech $X=[0,1]$ z $\sigma$-algebrą borelowską,
$\mu$ miarą liczącą, a $\nu$ ograniczeniem miary
Lebesgue'a do zbiorów borelowskich. Miara $\mu$
nie jest $\sigma$-skończona, gdyż każdy jej zbiór
skończonej miary jest skończony, a przeliczalna suma
takich zbiorów nie pokrywa nieprzeliczalnego $X$.
Ponieważ jedynym zbiorem $\mu$-zerowym jest $\varnothing$,
mamy $\nu\ll\mu$. Przypuśćmy, że
$\nu(E)=\int_Ef\,\mathrm d\mu(x)$ dla pewnej
mierzalnej $f\ge0$. Dla każdego $x\in X$:
\[
0=\nu(\{x\})=\int_{\{x\}}f\,\mathrm d\mu(x)=f(x).
\]
Zatem $f=0$ wszędzie, wbrew
$\nu(X)=1$. Jest to nawet przykład z miarą $\mu$
\emph{półskończoną}: każdy zbiór nieskończonej miary
zawiera niepusty zbiór skończonej dodatniej miary.
\end{przyklad}

\begin{przyklad}[Dlaczego $\sigma$-skończoność samej $\mu$ nie wystarcza?]
Na $X=(0,1)$ z $\mu=\lambda_1|_X$ definiujemy
$\nu(E)=0$ dla $\mu(E)=0$ oraz
$\nu(E)=\infty$ dla $\mu(E)>0$.
Jest to miara nieujemna, bezwzględnie ciągła względem
$\mu$, lecz nie jest $\sigma$-skończona: każdy zbiór
skończonej miary $\nu$ jest $\mu$-zerowy, a przeliczalna
suma takich zbiorów nie może pokryć $X$.

Nie istnieje gęstość $f:X\to[0,\infty)$, skończona
$\mu$-prawie wszędzie, reprezentująca $\nu$.
Istotnie, gdyby istniała, to dla pewnego $m\ge1$
zbiór $E_m=\{x:f(x)\le m\}$ miałby dodatnią miarę
$\mu$, ponieważ $\bigcup_{m\ge1}E_m$ ma pełną miarę.
Wtedy
$\int_{E_m}f(x)\,\mathrm d\mu(x)\le m\mu(E_m)<\infty$,
podczas gdy $\nu(E_m)=\infty$. Jest to sprzeczność.
Dopuszczenie gęstości przyjmującej wartość $\infty$
prowadzi do innego sformułowania i nie daje wniosku
o skończoności pochodnej prawie wszędzie.
\end{przyklad}

\section{Pochodna miary ze znakiem i miary zespolonej}
\begin{twierdzenie}[Radon--Nikodym dla miar zespolonych]
\label{thm:RN-complex-18}
Niech $\mu$ będzie nieujemną miarą $\sigma$-skończoną,
a $\eta$ skończoną miarą zespoloną na $(X,\mathcal F)$.
Jeżeli $\eta\ll\mu$, to istnieje dokładnie jedna
(z dokładnością do równości $\mu$-prawie wszędzie)
funkcja $f\in L^1(\mu;\mathbb C)$, dla której
\begin{equation}\label{eq:RN-complex-18}
\eta(E)=\int_E f(x)\,\mathrm d\mu(x)
\qquad(E\in\mathcal F).
\end{equation}
Ponadto $|\eta|(E)=\int_E|f|\,\mathrm d\mu(x)$,
a w szczególności $\|\eta\|_{\mathrm{TV}}=\|f\|_1$.
Dla rzeczywistej miary ze znakiem $\eta$ funkcję $f$
można wybrać rzeczywistą.
\end{twierdzenie}
\begin{proof}
Jeśli $\eta$ jest rzeczywistą skończoną miarą ze znakiem,
rozważamy jej rozkład Jordana
$\eta=\eta^+-\eta^-$ z rozdziału~\ref{ch:miary-ze-znakiem}.
Ponieważ $\eta\ll\mu$, dla każdego mierzalnego
zbioru $E$ o $\mu(E)=0$ również wszystkie podzbiory
mierzalne $A\subseteq E$ spełniają $\eta(A)=0$.
Z wzorów na części Jordana wynika
$\eta^+(E)=\eta^-(E)=0$. Obie miary dodatnie są
skończone, a więc $\sigma$-skończone, i są
bezwzględnie ciągłe względem $\mu$.
Stosujemy twierdzenie~\ref{thm:RN-sigfin-18} do każdej
z nich, otrzymując $f_+,f_-\in L^1(\mu)$.
Wtedy $f=f_+-f_-$ jest żądaną reprezentacją.
Dla miary zespolonej stosujemy ten argument do
$\operatorname{Re}\eta$ i $\operatorname{Im}\eta$
i otrzymujemy $f=f_1+if_2\in L^1(\mu;\mathbb C)$.
Jednoznaczność wynika z jednoznaczności części
rzeczywistych i urojonych, a wzór na wariację
jest twierdzeniem~\ref{thm:density-variation-16}.
\end{proof}

\begin{wniosek}[Pełny rozkład miary na część z gęstością i osobliwą]
\label{cor:RN-lebesgue-sum-18}
Niech $\mu,\nu$ będą nieujemnymi miarami
$\sigma$-skończonymi na $(X,\mathcal F)$.
Istnieją dokładnie jedna (z dokładnością do
$\mu$-prawie wszędzie) funkcja mierzalna
$f:X\to[0,\infty)$ i dokładnie jedna miara
$\nu_{\mathrm{s}}\perp\mu$ takie, że
\begin{equation}\label{eq:RN-Lebesgue-decomp-18}
\nu(E)=\int_E f(x)\,\mathrm d\mu(x)+\nu_{\mathrm{s}}(E)
\qquad(E\in\mathcal F).
\end{equation}
Dla skończonej miary zespolonej $\eta$ i
$\sigma$-skończonej $\mu$ analogiczna reprezentacja ma
$f\in L^1(\mu;\mathbb C)$ i osobliwą miarę zespoloną.
\end{wniosek}
\begin{proof}
Twierdzenie~\ref{thm:lebesgue-decomp-sigfin-17}
daje jednoznaczny rozkład $\nu=\nu_{\mathrm{ac}}+
\nu_{\mathrm{s}}$. Obie części są $\sigma$-skończone,
a więc \ref{thm:RN-sigfin-18} daje
$\nu_{\mathrm{ac}}(E)=\int_Ef\,\mathrm d\mu(x)$.
Jednoznaczność wynika z jednoznaczności obu składników
i pochodnej. Wersję zespoloną otrzymujemy z
wniosku~\ref{cor:complex-lebesgue-decomp-17}
i twierdzenia~\ref{thm:RN-complex-18}.
\end{proof}

\begin{twierdzenie}[Rozkład biegunowy miary zespolonej]
\label{thm:polar-complex-measure-18}
Niech $\eta$ będzie skończoną miarą zespoloną
na $(X,\mathcal F)$, a $|\eta|$ jej wariacją
całkowitą. Istnieje funkcja mierzalna
$u:X\to\mathbb C$, określona jednoznacznie
$|\eta|$-prawie wszędzie, taka, że
\begin{equation}\label{eq:polar-complex-measure-18}
 |u(x)|=1\quad |\eta|\text{-prawie wszędzie},\qquad
 \eta(E)=\int_Eu(x)\,\mathrm d|\eta|(x)
 \quad(E\in\mathcal F).
\end{equation}
Jeśli $\eta$ jest rzeczywistą miarą ze znakiem, to
$u$ można wybrać o wartościach w $\{-1,1\}$.
\end{twierdzenie}
\begin{proof}
Z definicji wariacji wynika $\eta\ll|\eta|$.
Miara $|\eta|$ jest skończona, toteż twierdzenie
\ref{thm:RN-complex-18} daje $u\in L^1(|\eta|;\mathbb C)$
spełniającą drugą równość w~\eqref{eq:polar-complex-measure-18}.
Twierdzenie~\ref{thm:density-variation-16} zastosowane
do tej reprezentacji daje
$|\eta|(E)=\int_E|u(x)|\,\mathrm d|\eta|(x)$.
Z drugiej strony $|\eta|(E)=\int_E1\,\mathrm d|\eta|(x)$.
Jednoznaczność pochodnej Radona--Nikodyma implikuje
$|u|=1$ $|\eta|$-prawie wszędzie.
W przypadku rzeczywistym gęstość jest rzeczywista,
więc warunek $|u|=1$ daje $u\in\{-1,1\}$
prawie wszędzie. Ustalając wartość $1$ na pozostałym
zbiorze zerowym, otrzymujemy żądaną funkcję mierzalną.
\end{proof}

\begin{przyklad}[Rozkład biegunowy jako znak gęstości]
Niech $\eta(E)=\int_E f(x)\,\mathrm d\mu(x)$,
gdzie $f\in L^1(\mu;\mathbb C)$. Wówczas
$|\eta|(E)=\int_E|f(x)|\,\mathrm d\mu(x)$,
a funkcję z poprzedniego twierdzenia można przyjąć
równą $u=f/|f|$ na $\{f\neq0\}$ oraz $u=1$
na $\{f=0\}$. W szczególności dla rzeczywistego $f$
jest to odpowiednik znaku funkcji.
\end{przyklad}

\section{Reguły rachunku dla pochodnych miar}
\begin{twierdzenie}[Reguła łańcuchowa]
\label{thm:RN-chain-18}
Niech $\mu,\nu,\rho$ będą nieujemnymi miarami
$\sigma$-skończonymi na $(X,\mathcal F)$
i niech $\rho\ll\nu\ll\mu$.
Oznaczmy $f=\mathrm d\nu/\mathrm d\mu$
i $g=\mathrm d\rho/\mathrm d\nu$.
Wtedy $\rho\ll\mu$ i
\begin{equation}\label{eq:RN-chain-18}
\frac{\mathrm d\rho}{\mathrm d\mu}
=g f\qquad\mu\text{-prawie wszędzie},
\end{equation}
przy czym na zbiorze $\{f=0\}$ przyjmujemy $gf=0$
niezależnie od wartości reprezentanta $g$.
\end{twierdzenie}
\begin{proof}
Bezwzględna ciągłość $\rho\ll\mu$ wynika
z przechodniości implikacji dla zbiorów zerowych.
Z lematu~\ref{lem:integration-density-18}, dla każdego
$E\in\mathcal F$,
\[
\rho(E)=\int_Eg\,\mathrm d\nu(x)
=\int_Eg f\,\mathrm d\mu(x).
\]
Ponieważ zmiana $g$ na zbiorze $\nu$-zerowym nie zmienia
całki względem $\nu$, a $\nu(\{f=0\})=0$,
przyjęta konwencja jest spójna. Jednoznaczność
pochodnej daje \eqref{eq:RN-chain-18}.
\end{proof}

\begin{twierdzenie}[Pochodna odwrotna]
\label{thm:RN-inverse-18}
Niech $\mu,\nu$ będą nieujemnymi miarami
$\sigma$-skończonymi i $\mu\sim\nu$.
Jeżeli $f=\mathrm d\nu/\mathrm d\mu$, to
$f>0$ $\mu$-prawie wszędzie i
\begin{equation}\label{eq:RN-inverse-18}
\frac{\mathrm d\mu}{\mathrm d\nu}=\frac1f
\qquad\nu\text{-prawie wszędzie}.
\end{equation}
\end{twierdzenie}
\begin{proof}
Z równoważności $\mu\ll\nu$ i wzoru
$\nu(\{f=0\})=\int_{\{f=0\}}f\,\mathrm d\mu(x)=0$
wynika $\mu(\{f=0\})=0$. Zatem $1/f$ jest
zdefiniowana i skończona prawie wszędzie.
Dla każdego $E\in\mathcal F$ lemat
\ref{lem:integration-density-18} daje
\[
\int_E\frac1f\,\mathrm d\nu(x)
=\int_E\frac1f f\,\mathrm d\mu(x)=\mu(E).
\]
Jednoznaczność pochodnej kończy dowód.
\end{proof}

\begin{wniosek}[Różnica i suma miar z gęstościami]
\label{cor:RN-linearity-18}
Niech $\mu$ będzie miarą $\sigma$-skończoną,
a $\eta_1,\eta_2$ skończonymi miarami zespolonymi,
bezwzględnie ciągłymi względem $\mu$.
Dla $a,b\in\mathbb C$ zachodzi
\[
\frac{\mathrm d(a\eta_1+b\eta_2)}{\mathrm d\mu}
=a\frac{\mathrm d\eta_1}{\mathrm d\mu}
+b\frac{\mathrm d\eta_2}{\mathrm d\mu}
\quad\mu\text{-prawie wszędzie}.
\]
\end{wniosek}
\begin{proof}
Wstawiamy reprezentacje całkowe z twierdzenia
\ref{thm:RN-complex-18}, korzystamy z liniowości całki,
a następnie z jednoznaczności gęstości.
\end{proof}

\section{Zastosowanie: warunkowa wartość oczekiwana}

\noindent\emph{Powiązanie tematyczne w Bibliotece matematycznej:} \atlasref{P1-10::ch:condexp}{P1-10: warunkowa wartość oczekiwana i rzut ortogonalny}.

Przedstawimy zastosowanie pokazujące, że twierdzenie
Radona--Nikodyma nie jest jedynie wynikiem o porównywaniu
miar. Pozwala ono skonstruować jeden z podstawowych
operatorów teorii prawdopodobieństwa.

\begin{definicja}[Warunkowa wartość oczekiwana]
\label{def:conditional-exp-18}
Niech $(\Omega,\mathcal F,\mathbb P)$ będzie przestrzenią
probabilistyczną, $\mathcal G\subseteq\mathcal F$
pod-$\sigma$-algebrą, a $Y\in L^1(\mathbb P;\mathbb R)$.
\emph{Warunkową wartością oczekiwaną} $Y$ względem
$\mathcal G$ nazywamy klasę funkcji
$Z\in L^1(\mathbb P;\mathbb R)$, która jest
$\mathcal G$-mierzalna i spełnia
\begin{equation}\label{eq:CE-def-18}
\int_A Z(\omega)\,\mathrm d\mathbb P(\omega)
=\int_A Y(\omega)\,\mathrm d\mathbb P(\omega)
\qquad(A\in\mathcal G).
\end{equation}
Piszemy $Z=\mathbb E[Y\mid\mathcal G]$.
Równość jest rozumiana z dokładnością do
$\mathbb P$-prawie wszystkich punktów. Ściślej, pracujemy
z klasami funkcji mierzalnych względem $\mathcal G$;
nie każda modyfikacja na zbiorze $\mathbb P$-zerowym musi
pozostać $\mathcal G$-mierzalna, gdy $\mathcal G$ nie jest zupełna.
\end{definicja}

\begin{twierdzenie}[Istnienie, jednoznaczność i własności podstawowe]
\label{thm:CE-existence-18}
Dla każdej przestrzeni probabilistycznej
$(\Omega,\mathcal F,\mathbb P)$, każdej
pod-$\sigma$-algebry $\mathcal G$ oraz $Y\in L^1(\mathbb P;\mathbb R)$
istnieje dokładnie jedna warunkowa wartość oczekiwana
$Z=\mathbb E[Y\mid\mathcal G]$ w sensie równości
$\mathbb P$-prawie wszędzie. Ponadto:
\begin{enumerate}[label=\textup{(\roman*)}]
\item odwzorowanie $Y\mapsto\mathbb E[Y\mid\mathcal G]$
jest liniowe i $\|Z\|_{L^1}\le\|Y\|_{L^1}$;
\item $Y\ge0$ prawie wszędzie implikuje $Z\ge0$
prawie wszędzie;
\item jeśli $Y$ jest $\mathcal G$-mierzalna,
to $Z=Y$ prawie wszędzie;
\item jeśli $\mathcal H\subseteq\mathcal G\subseteq\mathcal F$,
to $\mathbb E[\mathbb E[Y\mid\mathcal G]\mid\mathcal H]
=\mathbb E[Y\mid\mathcal H]$ prawie wszędzie.
\end{enumerate}
\end{twierdzenie}
\begin{proof}
Na przestrzeni mierzalnej $(\Omega,\mathcal G)$
określamy skończoną miarę ze znakiem
$\eta(A)=\int_A Y\,\mathrm d\mathbb P$.
Jest ona bezwzględnie ciągła względem
$\mu=\mathbb P|_{\mathcal G}$, a $\mu$ jest miarą
skończoną. Twierdzenie~\ref{thm:RN-complex-18}
(dla rzeczywistych miar) daje $\mathcal G$-mierzalną
$Z\in L^1(\mu)$, która spełnia
\eqref{eq:CE-def-18}. Jednoznaczność wynika z
jednoznaczności pochodnej Radona--Nikodyma.

Liniowość wynika z liniowości całki po $A$ i
jednoznaczności. Jeśli $Y\ge0$, to $\eta$ jest
miarą nieujemną, a więc jej pochodna $Z$ jest
nieujemna. Jeśli $Y$ jest $\mathcal G$-mierzalna,
sama spełnia \eqref{eq:CE-def-18}, co daje (iii).
Aby dowieść oszacowania $L^1$, korzystamy z
$|\eta|(\Omega)=\int_\Omega|Z|\,\mathrm d\mathbb P$
oraz z faktu, że dla każdego skończonego podziału
$\Omega=\bigsqcup_jA_j$, $A_j\in\mathcal G$,
\[
\sum_j|\eta(A_j)|
\le\sum_j\int_{A_j}|Y|\,\mathrm d\mathbb P
=\int_\Omega|Y|\,\mathrm d\mathbb P.
\]
Biorąc supremum po podziałach otrzymujemy (i).
Dla (iv) niech $A\in\mathcal H$; wówczas także
$A\in\mathcal G$, więc
\[
\int_A\mathbb E[Y\mid\mathcal G] \,\mathrm d\mathbb P
=\int_A Y\,\mathrm d\mathbb P.
\]
Zatem obie strony tożsamości (iv) spełniają
warunek definiujący wartość oczekiwaną względem
$\mathcal H$; jednoznaczność kończy dowód.
\end{proof}

\begin{twierdzenie}[Przypadek $L^2$: rzut ortogonalny]
\label{thm:CE-orthogonal-18}
W założeniach definicji~\ref{def:conditional-exp-18},
jeśli $Y\in L^2(\mathbb P;\mathbb R)$, to
$Z=\mathbb E[Y\mid\mathcal G]\in L^2(\mathbb P)$,
$\|Z\|_2\le\|Y\|_2$ i
\[
\int_\Omega(Y-Z)W\,\mathrm d\mathbb P=0
\qquad\text{dla każdego }W\in L^2(\Omega,\mathcal G,\mathbb P).
\]
Oznacza to, że $Z$ jest rzutem ortogonalnym $Y$
na domkniętą podprzestrzeń funkcji $\mathcal G$-mierzalnych.
\end{twierdzenie}
\begin{proof}
Nie wiemy jeszcze, czy $Z\in L^2$.
Niech $Z_n=\max(-n,\min(Z,n))$; jest to ograniczona
funkcja $\mathcal G$-mierzalna, więc z definicji
warunkowej wartości oczekiwanej (najpierw dla funkcji
prostych, potem przez zbieżność zmajoryzowaną) mamy
$\int_\Omega ZZ_n\,\mathrm d\mathbb P
=\int_\Omega YZ_n\,\mathrm d\mathbb P$.
Ponieważ $ZZ_n\ge Z_n^2$ punktowo,
\[
\|Z_n\|_2^2\le\int_\Omega ZZ_n\,\mathrm d\mathbb P
\le\|Y\|_2\|Z_n\|_2.
\]
Zatem $\|Z_n\|_2\le\|Y\|_2$.
Ponieważ $Z_n^2\uparrow Z^2$, Beppo Levi daje
$\|Z\|_2\le\|Y\|_2$. Definicyjna równość
$\int_A(Y-Z)\,\mathrm d\mathbb P=0$ dla $A\in\mathcal G$
rozciąga się liniowo na funkcje proste i przez
zbieżność zmajoryzowaną na ograniczone funkcje
$\mathcal G$-mierzalne. Trunkacje dowolnego
$W\in L^2(\mathcal G)$ zbliżają się do niego w $L^2$;
Cauchy--Schwarz pozwala przejść do granicy.
Podprzestrzeń $L^2(\mathcal G)$ jest domknięta:
jeśli $W_n$ są $\mathcal G$-mierzalne i zbieżne w $L^2$
do $W$, to wybieramy podciąg $(W_{n_k})$ taki, że
$\|W_{n_k}-W\|_2<2^{-2k}$. Z nierówności
Czebyszewa i lematu Borela--Cantellego wynika
$W_{n_k}\to W$ prawie wszędzie względem $\mathbb P$.
Zbiór $D$ punktów, w których ciąg $(W_{n_k})$
ma granicę skończoną, jest $\mathcal G$-mierzalny,
gdyż warunek Cauchy'ego i skończoność wartości
opisuje się przeliczalnymi działaniami mierzalnymi.
Funkcja $\widetilde W=\lim_kW_{n_k}$ na $D$
i $\widetilde W=0$ poza $D$ jest
$\mathcal G$-mierzalna i równa $W$ prawie wszędzie.
To konstruuje przedstawiciela w $L^2(\mathcal G)$,
bez zakładania zupełności $\mathcal G$.
Równanie ortogonalności identyfikuje $Z$ z rzutem.
\end{proof}

\section{Przykłady obliczania pochodnych}
\begin{przyklad}[Nieskończona gęstość przy punkcie brzegowym]
Na $X=(0,1)$ z miarą Lebesgue'a określmy
$\nu(E)=\int_E x^{-1/2}\,\mathrm d\lambda_1(x)$.
Wtedy $\nu\ll\lambda_1$, a
$\mathrm d\nu/\mathrm d\lambda_1=x^{-1/2}$.
Choć gęstość rośnie bez ograniczenia przy $0$,
jest całkowalna i $\nu(X)=2$.
\end{przyklad}

\begin{przyklad}[Miara osobliwa nie ma pochodnej reprezentującej masę]
Niech $X=[0,1]$ i $\mu=\lambda_1|_X$ na
$\mathcal B(X)$. Dla miary Cantora $\kappa$ z przykładu
\ref{ex:cantor-measure-17} zachodzi $\kappa\perp\mu$.
Nie istnieje $f\in L^1(\mu)$ spełniająca
$\kappa(E)=\int_Ef(x)\,\mathrm d\mu(x)$ dla wszystkich
$E\in\mathcal B(X)$, bo przy $E=C$ prawa strona
jest równa zero, a $\kappa(C)=1$.
Fakt, że pochodna klasycznej funkcji Cantora wynosi
zero prawie wszędzie, nie oznacza, że jej miara
Stieltjesa jest zerowa. Jest ona osobliwa względem
miary Lebesgue'a.
\end{przyklad}

\section{Rozkład miary Lebesgue'a--Stieltjesa}

\begin{twierdzenie}[Kryterium oparte na rozkładzie miary]\label{it:thm:stieltjes-decomposition}
Niech $\mu_\alpha$ będzie skończoną miarą Lebesgue'a--Stieltjesa ograniczoną do $(a,b]$. Istnieje rozkład
\[
\mu_\alpha=w\lambda_1+\mu_{\mathrm{sc}}+\sum_{k\in J}c_k\delta_{t_k},
\]
gdzie wszystkie miary rozpatrujemy na $(a,b]$, $w\ge0$ jest mierzalna, $\mu_{\mathrm{sc}}\perp\lambda_1$ jest bezatomowa, $J$ jest zbiorem co najwyżej przeliczalnym, $t_k$ są różne, $c_k>0$ i $\sum_{k\in J}c_k<\infty$. Dla każdej funkcji borelowskiej $f:(a,b]\to\mathbb{R}$ zachodzi
\[
\int_{(a,b]}|f|\,\mathrm d\mu_\alpha
=\int_{(a,b]}|f|w\,\mathrm d\lambda_1+
\int_{(a,b]}|f|\,\mathrm d\mu_{\mathrm{sc}}+
\sum_{k\in J}c_k|f(t_k)|.
\]
Funkcja jest bezwzględnie całkowalna względem $\mu_\alpha$ wtedy i tylko wtedy, gdy prawa strona jest skończona. W tym przypadku
\[
 \int_{(a,b]}f\,\mathrm d\mu_\alpha
 =\int_{(a,b]}fw\,\mathrm d\lambda_1+
 \int_{(a,b]}f\,\mathrm d\mu_{\mathrm{sc}}+
 \sum_{k\in J}c_k f(t_k),
\]
a ostatni szereg jest bezwzględnie zbieżny.
\end{twierdzenie}
\begin{proof}
Obie miary $\mu_\alpha$ i $\lambda_1|_{(a,b]}$ są skończone, więc twierdzenie o rozkładzie Lebesgue'a daje część bezwzględnie ciągłą i osobliwą. Twierdzenie \ref{thm:RN-sigfin-18} Radona--Nikodyma daje nieujemną gęstość $w$ pierwszej części względem $\lambda_1$. Drugą część przedłużamy przez zero na $\mathbb{R}$ i stosujemy twierdzenie \atlasref{ME-17::thm:atomic-cont-singular-17}{ME-17: rozkład części osobliwej}. Masa atomów sumuje się do masy części atomowej, nie większej od $\mu_\alpha((a,b])$.
Ponieważ część bezwzględnie ciągła nie ma mas punktowych,
każdy współczynnik $c_k$ jest równy
$\mu_\alpha(\{t_k\})$. Dla generatora prawostronnie
ciągłego oznacza to $c_k=\alpha(t_k)-\alpha(t_k-)$.
Jednoznaczność rozkładu Lebesgue'a i gęstości
Radona--Nikodyma daje jednoznaczność $w$ prawie wszędzie
względem $\lambda_1$ oraz części osobliwej.
Masy singletonów wyznaczają jej część atomową,
a odjęcie tej części wyznacza $\mu_{\mathrm{sc}}$.
Dla nieujemnej borelowskiej funkcji prostej równość całek względem sumy miar wynika z definicji całki i przeliczalnej addytywności sum nieujemnych. Dla dowolnej nieujemnej funkcji mierzalnej przybliżamy ją od dołu rosnącym ciągiem funkcji prostych i stosujemy twierdzenie \atlasref{ME-08::thm:mct}{ME-08: twierdzenie Beppo Leviego} o zbieżności monotonicznej do każdej miary oraz do sumy po $J$ (miara licząca). Lemat \ref{lem:integration-density-18} identyfikuje całkę względem $w\lambda_1$, a całka względem $\delta_{t_k}$ jest wartością w $t_k$. Stosujemy otrzymaną równość do $|f|$. Jeśli prawa strona jest skończona, stosujemy ją również osobno do $f^+$ i $f^-$ i odejmujemy skończone wartości. Z oszacowania dla $|f|$ wynika bezwzględna zbieżność szeregu z wartościami $f(t_k)$.
\end{proof}

\begin{twierdzenie}[Rozkład funkcji generującej miarę Stieltjesa]
\label{it:thm:stieltjes-generator-decomposition}
Niech $a<b$, a $g\in\operatorname{BV}[a,b]$ będzie prawostronnie ciągła
na $[a,b)$. Istnieją jednoznacznie określone funkcje
$g_{\mathrm{ac}},g_{\mathrm{sc}},g_{\mathrm{at}}\in\operatorname{BV}[a,b]$,
zerujące się w $a$, takie że
\[
 g(t)=g(a)+g_{\mathrm{ac}}(t)+g_{\mathrm{sc}}(t)+g_{\mathrm{at}}(t),
\]
gdzie
\[
 g_{\mathrm{ac}}(t)=\int_a^t w\,\mathrm d\lambda_1,\qquad
 g_{\mathrm{at}}(t)=\sum_{a<c\le t}d_c,
 \qquad w\inL^1[a,b],\quad \sum_c|d_c|<\infty.
\]
Funkcja $g_{\mathrm{sc}}$ jest ciągła, o ograniczonym
wahaniu, a jej pochodna znika prawie wszędzie.
Ponadto $w=g'$ prawie wszędzie i $d_c=g(c)-g(c-)$.
Sumy po $c$ obejmują punkty z $(a,b]$ o niezerowym
współczynniku; jest ich co najwyżej przeliczalnie wiele.
W zapisie miary $g'\lambda_1$ oraz we wzorach całkowych
wybieramy skończonego borelowskiego reprezentanta
pochodnej, określonej prawie wszędzie.
Zachodzi dokładna równość
\[
 V_a^b(g)=\int_a^b|g'|\,\mathrm d\lambda_1+
 V_a^b(g_{\mathrm{sc}})+\sum_{a<c\le b}|d_c|.
\]
Jeżeli $g$ jest niemalejąca, wszystkie trzy składniki
są niemalejące.
\end{twierdzenie}
\begin{proof}
Niech $\nu=\nu_g$ będzie miarą z twierdzenia
\atlasref{ME-16::it:thm:stieltjes-signed-bv}{ME-16}. Do dodatnich części
$\nu^+,\nu^-$ jej rozkładu Jordana stosujemy rozkład
z twierdzenia \ref{it:thm:stieltjes-decomposition}:
każda z tych miar jest miarą Stieltjesa dzięki twierdzeniu
\atlasref{ME-04::it:thm:stieltjes-construction}{ME-04} (przedłużamy ją przez
zero poza $(a,b]$). Odejmując odpowiadające sobie części,
otrzymujemy
\[
 \nu=w\lambda_1+\sigma+\sum_{c\in J}d_c\delta_c,
\]
gdzie $\sigma$ jest skończoną miarą ze znakiem osobliwą
względem $\lambda_1$ i jej wahanie jest bezatomowe.
Istotnie, wahanie różnicy dwóch dodatnich miar jest
ograniczone przez ich sumę, która w tym przypadku
jest osobliwa i bezatomowa. Na prostej bezatomowość
wynika z zerowych mas wszystkich singletonów
(twierdzenie \atlasref{ME-17::thm:atomic-cont-singular-17}{ME-17}).
Zbiór $J$ jest co najwyżej przeliczalny, a
$d_c=\nu(\{c\})=g(c)-g(c-)$, ponieważ pozostałe
części nie mają atomów. Stąd
$\sum_{c\in J}|d_c|\le|\nu|((a,b])$.

Wybieramy borelowski zbiór zerowy $S$, na którym skupia się
$|\sigma|$, i dołączamy $J$ do $S$. Trzy miary ze znakiem
są skupione odpowiednio na rozłącznych zbiorach
$(a,b]\setminus S$, $S\setminus J$, $J$.
Ich wahania zatem sumują się do $|\nu|$:
na każdym z tych trzech zbiorów ograniczenie $\nu$
zgadza się z odpowiednim składnikiem, a wahanie
ograniczenia jest ograniczeniem wahania.
Jednoznaczność części atomowej wynika z wartości na
singletonach. Po jej odjęciu, gdyby istniały dwa rozkłady
$\nu-\sum_c d_c\delta_c=w_i\lambda_1+\sigma_i$,
$i=1,2$, to $(w_1-w_2)\lambda_1=\sigma_2-\sigma_1$.
Miara po lewej jest bezwzględnie ciągła, a po prawej
skupiona na sumie dwóch borelowskich zbiorów zerowych.
Miara ze znakiem o obu tych własnościach jest zerowa:
każdy zbiór borelowski rozcinamy na część w tym zbiorze
zerowym i część poza nim. Stąd $\sigma_1=\sigma_2$ oraz $w_1=w_2$ prawie
wszędzie. Ostatnia implikacja wynika z jednoznaczności
gęstości Radona--Nikodyma lub bezpośrednio przez
całkowanie różnicy na zbiorach, gdzie jest ona dodatnia
i gdzie jest ujemna. Gdy gęstości są jedynie mierzalne
w sensie Lebesgue'a, wybieramy borelowskie wersje
zgodne z nimi prawie wszędzie.

Definiujemy składniki funkcji przez masy $(a,t]$
odpowiednich miar; w $a$ przyjmujemy zero.
Twierdzenie \atlasref{ME-16::it:thm:stieltjes-signed-bv}{ME-16} zapewnia
prawostronną ciągłość, ograniczone wahanie i jego dokładną
wartość dla każdego składnika. Dla $\sigma$ dodatkowo
$g_{\mathrm{sc}}(t)-g_{\mathrm{sc}}(t-)=\sigma(\{t\})=0$,
więc $g_{\mathrm{sc}}$ jest ciągła także lewostronnie.
Twierdzenie \atlasref{ME-25::thm:cdf-derivative-25}{ME-25: pochodna funkcji rozkładu}, zastosowane
do dodatnich części miar przedłużonych przez zero na
$\mathbb R$, daje $g'=w$ prawie wszędzie oraz zerową pochodną
części osobliwej i atomowej prawie wszędzie.
W szczególności $g_{\mathrm{sc}}'=0$ prawie wszędzie.
W powyższej konstrukcji $w$ jest różnicą borelowskich
gęstości Radona--Nikodyma dla miar na borelowskich
podzbiorach $(a,b]$, zatem można użyć $w$ jako
reprezentanta $g'$ (zmieniając wartości nieskończone
na zbiorze zerowym oraz wybierając wartość w $a$).
Zapisy z $g'$ w dalszych wzorach nie zależą od
wyboru takiego reprezentanta, ponieważ dotyczą
wyłącznie całkowania względem $\lambda_1$.

W celu sprawdzenia jednoznaczności zwróćmy uwagę,
że każda funkcja postaci
$g_{\mathrm{at}}(t)=\sum_{a<c\le t}d_c$ jest prawostronnie
ciągła. Jest bowiem jednostajną granicą skończonych sum
funkcji $d_c\chi_{[c,b]}$: norma supremum ogona jest
ograniczona przez ogon zbieżnego szeregu $\sum_c|d_c|$.
Każda skończona suma jest prawostronnie ciągła.
Pozostałe dwa składniki są ciągłe z założeń.
Wszystkie spełniają zatem hipotezy twierdzenia
\atlasref{ME-16::it:thm:stieltjes-signed-bv}{ME-16}.
Jednoznaczność funkcji wynika z jednoznaczności rozkładu
miary: bezwzględnie ciągły składnik generuje miarę
$w\lambda_1$, ciągły składnik o zerowej pochodnej
generuje miarę osobliwą bezatomową: ciągłość wyklucza
atomy na mocy twierdzenia \atlasref{ME-16::it:thm:stieltjes-signed-bv}{ME-16},
a twierdzenie \atlasref{ME-25::thm:cdf-derivative-25}{ME-25: pochodna funkcji rozkładu} zastosowane
do obu części Jordana daje gęstość części bezwzględnie
ciągłej równą różnicy ich gęstości, czyli pochodnej
funkcji generującej, tutaj zerowej prawie wszędzie.
Składnik skokowy generuje miarę atomową, gdyż jego
przyrost na $(s,t]$ jest sumą współczynników
$d_c$ dla $s<c\le t$ i szereg jest bezwzględnie zbieżny.
Jednoznaczność miary na takich przedziałach kończy
identyfikację.
Normalizacja w $a$ usuwa dowolność stałych.
Doprecyzujmy przejście od rozkładu miar do równości
wahań. Dla borelowskiego $B\subset(a,b]$ zachodzi
\[
 |w\lambda_1|(B)=\int_B|w|\,\mathrm d\lambda_1.
\]
Istotnie, dla dowolnego skończonego mierzalnego podziału
$B$ suma modułów całek $w$ nie przekracza całki $|w|$.
Podział $B$ na części, na których $w\ge0$ i $w<0$,
daje równość. Podobnie dla miary atomowej
$\tau=\sum_c d_c\delta_c$ mamy
$|\tau|(B)=\sum_{c\in B}|d_c|$: nierówność górna wynika
z nierówności trójkąta dla bezwzględnie zbieżnych
szeregów, a dolna z podziału na zbiory punktów
o dodatnich i ujemnych współczynnikach oraz ich
dopełnienie. Z rozłączności zbiorów skupienia
trzech składników wynika zatem
\[
 |\nu|(B)=\int_B|w|\,\mathrm d\lambda_1
           +|\sigma|(B)+\sum_{c\in B}|d_c|.
\]
Dla $B=(a,b]$ stosujemy twierdzenie
\atlasref{ME-16::it:thm:stieltjes-signed-bv}{ME-16} do $g$ i $g_{\mathrm{sc}}$,
a następnie $w=g'$ prawie wszędzie. Otrzymujemy
dokładnie żądaną równość wahań funkcji. Dla niemalejącego $g$ miara $\nu$ jest
nieujemna i jej trzy składniki również są nieujemne.
\end{proof}

\begin{wniosek}[Kryterium całkowalności względem miary Stieltjesa ze znakiem]
\label{it:cor:stieltjes-signed-integrability}
W hipotezach twierdzenia \ref{it:thm:stieltjes-generator-decomposition}
niech $\nu_g=g'\lambda_1+\sigma+\sum_c d_c\delta_c$
będzie otrzymanym rozkładem miary.
Dla każdej rzeczywistej lub zespolonej funkcji borelowskiej
$f$ na $(a,b]$ zachodzi
\[
 \int_{(a,b]}|f|\,\mathrm d|\nu_g|
 =\int_{(a,b]}|f|\,|g'|\,\mathrm d\lambda_1
  +\int_{(a,b]}|f|\,\mathrm d|\sigma|
  +\sum_c|d_c|\,|f(c)|.
\]
Funkcja $f$ jest całkowalna względem $\nu_g$ wtedy i tylko
wtedy, gdy prawa strona jest skończona. Wówczas
\[
 \int_{(a,b]}f\,\mathrm d\nu_g
 =\int_{(a,b]}fg'\,\mathrm d\lambda_1
  +\int_{(a,b]}f\,\mathrm d\sigma+\sum_c d_cf(c),
\]
przy czym szereg jest bezwzględnie zbieżny.
\end{wniosek}
\begin{proof}
Równość miar dodatnich
$|\nu_g|=|g'|\lambda_1+|\sigma|+\sum_c|d_c|\delta_c$
wykazano w dowodzie poprzedniego twierdzenia.
Argument dla całek nieujemnych z twierdzenia
\ref{it:thm:stieltjes-decomposition} daje pierwszą równość,
w szczególności kryterium całkowalności. Jeżeli prawa
strona jest skończona, to wszystkie trzy całki względem
składników miary ze znakiem istnieją, a szereg wartości
jest bezwzględnie zbieżny. Dla rzeczywistych funkcji
prostych równość wynika z addytywności rozkładu miary.
Dla funkcji nieujemnych całkowalnych przechodzimy
do granicy rosnących funkcji prostych, stosując
zbieżność monotoniczną osobno do dodatnich i ujemnych
części każdej miary. Wszystkie otrzymane granice są
skończone dzięki pierwszej równości. Rozkład funkcji
rzeczywistej na część dodatnią i ujemną oraz rozkład
funkcji zespolonej na część rzeczywistą i urojoną
kończą dowód.
\end{proof}

\section{Mapa implikacji}
Przy $\sigma$-skończoności $\mu,\nu$ uzyskaliśmy
następujący ciąg rezultatów: rozkład Lebesgue'a
$\nu=\nu_{\mathrm{ac}}+\nu_{\mathrm{s}}$,
reprezentację $\nu_{\mathrm{ac}}=f\mu$,
jednoznaczność gęstości, reguły rachunkowe i
konstrukcję warunkowej wartości oczekiwanej.
Przykład~\ref{ex:RN-failure-nonsigma-18} ostrzega,
że sama implikacja między zbiorami zerowymi
nie zastępuje hipotez twierdzenia.
